% ============================================
% Chapter 10: Conclusion
% ============================================

\chapter{Conclusion}

\section{Summary}

This report presented the design, implementation, and evaluation of the AEGIS Entropy system --- an AI-based intrusion detection and defense platform for port scan and network reconnaissance attacks. The project followed a rigorous methodology inspired by the CRISP-DM framework, encompassing data understanding, preparation, modeling, evaluation, and deployment.

\section{Key Contributions}

\begin{enumerate}
    \item \textbf{ML Detection Pipeline}: A complete data pipeline built on the CICIDS2017 dataset, achieving F1-Scores of 99.920\% (Random Forest) and 99.923\% (XGBoost) on an 11-feature schema, far exceeding the project targets of F1 $\geq$ 88\%, Accuracy $\geq$ 90\%, and FPR $<$ 6\%.

    \item \textbf{Defense and Deception Subsystem}: A layered defense architecture comprising honeypot redirection, phantom network simulation, Moving Target Defense (port mutation), traffic shaping (NetfilterQueue packet injection), and dynamic IP blacklisting --- all orchestrated by a severity-based Active Defense Engine.

    \item \textbf{Capture Environment}: A controlled VirtualBox laboratory with Kali Linux (attacker) and Ubuntu (target) machines, producing 10 Nmap scan scenarios covering stealthy to aggressive reconnaissance behaviors.

    \item \textbf{System Integration}: A unified Bridge module connecting Nmap XML parsing, ML prediction, and dashboard display, with dual logging (local JSONL + HTTP POST) ensuring no data loss.

    \item \textbf{SOC Dashboard}: A Flask + Socket.IO dashboard providing real-time alert visualization, blocked IP management, and honeypot event tracking.
\end{enumerate}

\section{Limitations}

\begin{itemize}
    \item Isolation Forest performance is poor on CICIDS2017 due to the 55\% attack ratio (majority anomaly paradox); effective deployment requires training on benign-only traffic
    \item The Traffic Shaper requires root privileges and \texttt{libnetfilter-queue}, limiting deployment to Linux environments
    \item Port rotation introduces brief connection drops during rule updates
    \item Phantom network is limited to ICMP/ARP; sophisticated attackers can detect phantom hosts via TCP probes
\end{itemize}

\section{Future Work}

\begin{itemize}
    \item \textbf{Adaptive MTD}: Use reinforcement learning to optimize rotation intervals based on threat level
    \item \textbf{Honeypot Fingerprinting}: Deploy realistic service emulators (SSH, HTTP, MySQL) instead of simple TCP listeners
    \item \textbf{Real-time Capture Integration}: Connect Scapy-based live capture to the ML pipeline for continuous monitoring
    \item \textbf{Multi-class Classification}: Extend the ML models to distinguish between scan types (SYN, UDP, ACK, XMAS, Connect)
    \item \textbf{Distributed MTD}: Extend port rotation across multiple servers in a cluster
\end{itemize}
